\section*{Chapitre \ref{ch:b3:palaeomagnetism} --- Géomagnétisme et paléomagnétisme}

\begin{solution}{exo:b3:palaeomagnetism:1}
$\tan I = 2\tan\ang{30} = 1.1547$, donc $I = \ang{49.1}$. Inversement, $\tan\lambda
= \tfrac12\tan(\ang{-20}) = -0.1820$, donc $\lambda = \ang{-10.3}$, soit
\ang{10.3}\,S.
\end{solution}

\begin{solution}{exo:b3:palaeomagnetism:2}
{\sloppy $H = \sqrt{20000^2 + 2000^2} = \qty{20100}{nT}$~;
$F = \sqrt{20000^2 + 2000^2 + 40000^2} = \qty{44770}{nT}$. La déclinaison
magnétique est $\arctan(-2000/20000) = \ang{-5.7}$, \ang{5.7} à l'ouest du nord~;
l'inclinaison $\arctan(40000/20100) = \ang{63.3}$, vers le bas.\par}
\end{solution}

\begin{solution}{exo:b3:palaeomagnetism:3}
$B_0 = -g_1^0 = \qty{29350}{nT}$ à l'équateur, et $2B_0 = \qty{58700}{nT}$ aux
pôles. $m = 4\pi a^3B_0/\mu_0 = a^3B_0/\num{e-7} = \num{2.5862e20}\times
\num{2.935e-5}/\num{e-7} = \qty{7.59e22}{A.m^2}$, un peu moins que le moment
\qty{7.69e22}{A.m^2} du dipôle complet, qui ajoute $g_1^1$ et $h_1^1$.
\end{solution}

\begin{solution}{exo:b3:palaeomagnetism:4}
(a) Une \omterm{def:b3:palaeomagnetism:nrm}{aimantation thermorémanente}, qui enregistre le champ au moment où la lave
se refroidit en traversant les \omterm{def:b3:palaeomagnetism:blocking}{températures de blocage} de ses grains. (b) Une \omterm{def:b3:palaeomagnetism:nrm}{aimantation rémanente détritique},
qui enregistre le champ pendant la sédimentation du limon,
jusqu'à son verrouillage par compaction. (c) Une \omterm{def:b3:palaeomagnetism:nrm}{aimantation rémanente chimique}, qui enregistre le
champ pendant la croissance de l'hématite, lors de l'enfouissement précoce et donc un peu
après le dépôt. (d) Une rémanence visqueuse le long du champ actuel, qui s'ajoute
à l'\omterm{def:b3:palaeomagnetism:nrm}{aimantation thermorémanente} d'origine, portée par les grains à temps de relaxation le plus court~; elle
est éliminée par les premières étapes de la désaimantation.
\end{solution}

\begin{solution}{exo:b3:palaeomagnetism:5}
Les vecteurs unitaires $(\cos I\cos D,\ \cos I\sin D,\ \sin I)$ sont
\begin{gather*}
  (0.6545, 0.1391, 0.7431),\quad (0.6670, 0.2695, 0.6947),\\
  (0.5235, 0.1501, 0.8387),\quad (0.5985, 0.0629, 0.7986).
\end{gather*}
{\sloppy La somme vaut $(2.4435, 0.6216, 3.0751)$, de longueur $R = 3.9766$. Moyenne~:
$D = \arctan(0.6216/2.4435) = \ang{14.3}$, $I = \arcsin(3.0751/3.9766) =
\ang{50.7}$. $k = 3/(4 - 3.9766) = 128$. Avec $20^{1/3} - 1 = 1.7144$,
$\cos\alpha_{95} = 1 - (0.0234/3.9766)\times1.7144 = 0.98991$, donc
$\alpha_{95} = \ang{8.1}$~: avec quatre échantillons, le cône reste large, quelle que soit
la cohésion du groupe.\par}
\end{solution}

\begin{solution}{exo:b3:palaeomagnetism:6}
{\sloppy $\lambda = \arctan(\tfrac12\tan\ang{50}) = \ang{30.79}$, donc $p = \ang{59.21}$.
$\sin\lambda_p = \sin\ang{40}\cos\ang{59.21} + \cos\ang{40}\sin\ang{59.21}
\cos\ang{20} = 0.9474$, $\lambda_p = \ang{71.34}$\,N. $\sin\beta =
\sin\ang{59.21}\sin\ang{20}/\cos\ang{71.34} = 0.9184$, $\beta = \ang{66.66}$.
Comme $\cos p = 0.5119 < \sin\ang{40}\times0.9474 = 0.6090$, le second cas
s'applique~: $\varphi_p = 10 + 180 - 66.66 = \ang{123.34}$\,E. Le pôle est à
\ang{71.3}\,N, \ang{123.3}\,E.\par}
\end{solution}

\begin{solution}{exo:b3:palaeomagnetism:7}
(a) $(7.69 - 8.32)/8.32 = \qty{-7.6}{\percent}$ en 125 ans, soit
\qty{-6.1}{\percent} par siècle. (b) $29350/12.6 = 2329$, environ 2300 ans.
Ce n'est pas une prévision~: le taux lui-même change d'un modèle quinquennal au
suivant, le moment a été à la fois plus élevé et plus faible dans le passé sans
s'inverser, et le champ s'inverse à des intervalles irréguliers que nul taux ne
prédit.
\end{solution}

\begin{solution}{exo:b3:palaeomagnetism:8}
Entre 20 et \qty{300}{\celsius}, le vecteur a perdu $(6.4, 0.0, 13.6)$~: une
composante de \qty{15.0}{mA/m}, avec $D = \ang{0.0}$ et
$I = \arctan(13.6/6.4) = \ang{64.8}$. Au-dessus de \qty{300}{\celsius}, il a perdu $(25.8, 18.1, 24.6)$~:
\qty{40.0}{mA/m}, $D = \arctan(18.1/25.8) = \ang{35.1}$, $I =
\arctan(24.6/31.5) = \ang{38.0}$. La seconde, qui passe par l'origine, est la
ChRM~; la première, dirigée vers le nord et très inclinée, éliminée à basse
température, est une surimpression visqueuse le long du champ récent.
\end{solution}

\begin{solution}{exo:b3:palaeomagnetism:9}
$\lambda = \arctan(\tfrac12\tan(\ang{-35})) = \ang{-19.3}$. Pour $I = \ang{-40}$
et \ang{-30}~: \ang{-22.8} et \ang{-16.1}, un intervalle de \ang{6.7}.
$p = \ang{90} - \lambda = \ang{109.3}$, $\cos^2p = 0.1096$~:
$\mathrm dp = 2.5\times(1 + 0.329) = \ang{3.3}$ et $\mathrm dm =
5\times\sin\ang{109.3}/\cos\ang{35} = \ang{5.8}$.
\end{solution}

\begin{solution}{exo:b3:palaeomagnetism:10}
(a) $\tan\lambda = \tfrac12$, $\lambda = \ang{26.6}$. $F$ vaut $B_0$ à
l'équateur et $2B_0$ aux pôles~: un rapport de 2. (b) $\tan I = 2\tan\lambda$
donne $(1 + \tan^2I)\,\mathrm dI = 2(1 + \tan^2\lambda)\,\mathrm d\lambda$, donc
\[
  \frac{\mathrm dI}{\mathrm d\lambda} = \frac{2/\cos^2\lambda}{1 + 4\tan^2\lambda}
  = \frac{2}{\cos^2\lambda + 4\sin^2\lambda} = \frac{2}{1 + 3\sin^2\lambda}.
\]
Une erreur $\delta I$ donne $\delta\lambda = \tfrac12(1 + 3\sin^2\lambda)\,\delta I$,
de $\tfrac12\delta I$ à l'équateur, jusqu'à $2\,\delta I$ aux pôles~: la
plus grande erreur se situe près des pôles. (c) Près de l'équateur, l'inclinaison varie
deux fois plus vite que la latitude, si bien qu'une erreur donnée d'inclinaison coûte
deux fois moins de latitude~; près des pôles, l'inclinaison varie à peine avec la latitude.
\end{solution}

\begin{solution}{exo:b3:palaeomagnetism:11}
(a) $\cos p = \sin(\ang{-35})\sin\ang{70} + \cos\ang{35}\cos\ang{70}\cos\ang{150}
= -0.5390 - 0.2426 = -0.7816$, $p = \ang{141.41}$. Puis $\tan I_x = 2\cot p =
-2.506$, $I_x = \ang{-68.2}$, et $\lambda_x = \ang{90} - p = \ang{-51.4}$.
$\cos D_x = (\sin\ang{70} - \sin(\ang{-35})\cos p)/(\cos\ang{35}\sin p) =
0.9617$, et comme le pôle se trouve à l'est du méridien du site ($\sin\ang{150} >
0$), $D_x = \ang{15.9}$. (b) $\lambda_o = \arctan(\tfrac12\tan(\ang{-60})) =
\ang{-40.9}$. $R = 30.0 - 15.9 = \ang{14.1}$~: une rotation dans le sens horaire du bloc
autour d'un axe vertical. $P = \lambda_x - \lambda_o = -51.4 + 40.9 =
\ang{-10.5}$~: la formation se trouvait \ang{10.5} plus au nord que ne le
prévoit la plaque stable, si bien que le bloc s'est depuis déplacé d'environ \ang{10.5}, soit
\qty{1170}{km} vers le sud par rapport à la plaque.
\end{solution}

\begin{solution}{exo:b3:palaeomagnetism:12}
(a) Avec $\mathrm d\Omega = 2\pi\sin\theta'\,\mathrm d\theta'$,
\[
  \int_0^\theta\frac{\kappa e^{\kappa\cos\theta'}}{4\pi\sinh\kappa}\,2\pi\sin\theta'\,
  \mathrm d\theta' = \frac{e^\kappa - e^{\kappa\cos\theta}}{2\sinh\kappa}
  = \frac{e^\kappa - e^{\kappa\cos\theta}}{e^\kappa - e^{-\kappa}} .
\]
(b) Pour $\kappa$ grand, $e^{-\kappa}$ est négligeable~; en divisant par $e^\kappa$
on obtient $1 - e^{-\kappa(1 - \cos\theta)} \approx 1 - e^{-\kappa\theta^2/2}$. Elle
vaut $1 - e^{-1} = \qty{63}{\percent}$ lorsque $\kappa\theta^2/2 = 1$, $\theta =
\sqrt{2/\kappa}$ radians, soit $\ang{81}/\sqrt\kappa$. (c) $81/\sqrt{50} =
\ang{11.5}$~; la formule exacte donne \qty{63.1}{\percent}.
\end{solution}

\begin{solution}{pb:b3:palaeomagnetism:1}
\textbf{1.} L'hématite des couches rouges peut se former à partir d'oxydes et d'hydroxydes de fer
bien après le dépôt du sédiment, ce qui confère une \omterm{def:b3:palaeomagnetism:nrm}{aimantation rémanente chimique} de cet
âge plus récent. Un test de pli positif, si les couches ont été plissées peu après
leur dépôt, ou un test d'inversion positif sur des couches rapprochées de
deux polarités montrerait que l'aimantation est précoce.

\textbf{2.} $(\cos\ang{37}\cos\ang{50}, \cos\ang{37}\sin\ang{50}, \sin\ang{37})
= (0.5134, 0.6118, 0.6018)$.

\textbf{3.} $R = \sqrt{3.8943^2 + 2.6500^2 + 3.5717^2} = 5.9114$. $D =
\arctan(2.6500/3.8943) = \ang{34.2}$, $I = \arcsin(3.5717/5.9114) = \ang{37.2}$.

\textbf{4.} $k = 5/(6 - 5.9114) = 56.5$. Avec $20^{1/5} - 1 = 0.8206$,
$\cos\alpha_{95} = 1 - (0.0886/5.9114)\times0.8206 = 0.98771$, $\alpha_{95} =
\ang{9.0}$.

\textbf{5.} B~: $k = 6/(7 - 6.8979) = 58.8$~; $20^{1/6} - 1 = 0.6475$,
$\cos\alpha_{95} = 1 - (0.1021/6.8979)\times0.6475 = 0.99042$, $\alpha_{95} =
\ang{7.9}$. C~: $k = 7/(8 - 7.8707) = 54.1$~; $20^{1/7} - 1 = 0.5340$,
$\cos\alpha_{95} = 1 - (0.1293/7.8707)\times0.5340 = 0.99123$, $\alpha_{95} =
\ang{7.6}$.

\textbf{6.} B a le $k$ le plus élevé, c'est le groupe le plus serré~; C a le plus petit
$\alpha_{95}$, la moyenne la mieux définie, car elle repose sur le plus grand nombre de sites. Les trois
sont proches~: les trois formations sont bien regroupées.

\textbf{7.} $\lambda_A = \arctan(\tfrac12\tan\ang{37.17}) = \ang{20.76}$\,N,
$\lambda_B = \ang{32.15}$\,N, $\lambda_C = \ang{45.33}$\,N.

\textbf{8.} $p = 90 - 20.76 = \ang{69.24}$. $\sin\lambda_p =
\sin\ang{47}\cos\ang{69.24} + \cos\ang{47}\sin\ang{69.24}\cos\ang{34.23} =
0.2593 + 0.5272 = 0.7865$~: $\lambda_p = \ang{51.86}$\,N.

\textbf{9.} $\sin\beta = \sin\ang{69.24}\sin\ang{34.23}/\cos\ang{51.86} =
0.8516$, $\beta = \ang{58.40}$. Puisque $\cos p = 0.3545 < \sin\ang{47}\times
0.7865 = 0.5752$, le deuxième cas s'applique, $\cos\beta < 0$~:
$\varphi_p = 15.00 + 180 - 58.40 = \ang{136.60}$\,E.

\textbf{10.} $\cos^2p = 0.1257$~: $\mathrm dp = 4.50\times1.377 = \ang{6.2}$~;
$\mathrm dm = 8.99\times\sin\ang{69.24}/\cos\ang{37.17} = \ang{10.5}$.

\textbf{11.} De \ang{51.9}\,N à \qty{60}{Ma} jusqu'à \ang{66.5}\,N à
\qty{30}{Ma} puis \ang{87.1}\,N à \qty{4}{Ma}, sur pratiquement le même méridien
(\ang{136}--\ang{139}\,E)~: une trajectoire vers le pôle géographique, atteinte dans
son cône par la formation la plus jeune.

\textbf{12.} Le niveau de cendres, à \qty{22}{m}, se situe dans la deuxième zone à polarité normale,
19.7--\qty{24.6}{m}~; avec un âge proche de \qty{3.2}{Ma}, cette zone est C2An.2n
(3.207--3.116). Puis la zone 4.2--\qty{14.8}{m} est C2An.3n, la zone
29.2--\qty{54.0}{m} C2An.1n, la base à polarité inverse appartient au chron situé sous
C2An, et le sommet à polarité inverse au chron situé au-dessus.

\textbf{13.} $(19.7 - 4.2)/(3.596 - 3.207) = 15.5/0.389 = \qty{39.8}{m/Ma}$~;
$(54.0 - 24.6)/(3.116 - 2.581) = 29.4/0.535 = \qty{55.0}{m/Ma}$.

\textbf{14.} Base~: $3.596 + 4.2/39.8 = \qty{3.70}{Ma}$. Sommet~: $2.581 -
10.0/55.0 = \qty{2.40}{Ma}$.

\textbf{15.} Seules les limites d'inversion sont datées, espacées de quelques centaines de milliers
d'années~; entre elles, le taux est supposé constant. Un changement de subsidence,
d'apport sédimentaire ou de climat, ou une lacune de dépôt, le modifierait, et un
hiatus pourrait même effacer une courte zone et fausser la corrélation.

\textbf{16.} $\cos p = \sin\ang{47}\sin\ang{78} + \cos\ang{47}\cos\ang{78}
\cos\ang{175} = 0.7154 - 0.1413 = 0.5741$, $p = \ang{54.96}$. $\tan I_x =
2\cot p = 1.4023$, $I_x = \ang{54.5}$~; $\lambda_x = \ang{35.04}$\,N.

\textbf{17.} $\cos D_x = (\sin\ang{78} - \sin\ang{47}\times0.5741)/
(\cos\ang{47}\sin\ang{54.96}) = 0.99976$~; le pôle se situe à l'est du
méridien du site ($\sin\ang{175} > 0$), donc $D_x = \ang{1.3}$.

\textbf{18.} \omterm{def:b3:palaeomagnetism:vgp}{Paléolatitudes} attendues à partir de $I_x$ = \ang{54.5}, \ang{60.0} et
\ang{63.6}~: \ang{35.03}, \ang{40.89} et \ang{45.21}. $P_A = 35.03 - 20.76 =
\ang{14.3}$, $P_B = 40.89 - 32.15 = \ang{8.7}$, $P_C = 45.21 - 45.33 =
\ang{-0.1}$.

\textbf{19.} $R_A = 34.2 - 1.3 = \ang{32.9}$, $R_B = 23.9 - 2.9 = \ang{21.0}$,
$R_C = 3.4 - 2.5 = \ang{0.9}$, toutes dans le sens horaire.

\textbf{20.} $\Delta D_o = \arcsin(\sin\ang{8.99}/\cos\ang{37.17}) = \ang{11.3}$~;
$\Delta D_x = \arcsin(\sin\ang{3.0}/\cos\ang{54.5}) = \ang{5.2}$~; $\Delta R_A =
\sqrt{11.3^2 + 5.2^2} = \ang{12.4}$.

\textbf{21.} $R_A = \ang{32.9} \pm \ang{12.4}$ et $R_B = \ang{21.0} \pm
\ang{13.8}$ sont significatifs~; $R_C = \ang{0.9} \pm \ang{17.9}$ ne l'est pas.

\textbf{22.} Il y a \qty{60}{Ma}, Kesta se trouvait environ \ang{14} de latitude
au sud de la position que lui assigne le continent stable, et avait pivoté d'environ \ang{33}
dans le sens antihoraire par rapport à son orientation actuelle. À \qty{30}{Ma}, il s'était
rapproché d'environ \ang{6} de ces \ang{14} vers le nord, et avait pivoté de \ang{12} sur les
\ang{33} initiaux dans le sens horaire~; à \qty{4}{Ma}, il avait atteint sa place actuelle sur le continent,
sans rotation restante. Il a dérivé vers le nord le long de la marge du continent, en pivotant
dans le sens horaire au fil de sa progression, et était en place dès \qty{4}{Ma}.

\textbf{23.} $47.00 - 20.76 = \ang{26.24}$, soit $26.24\times111.19 =
\qty{2918}{km}$. Cela inclut le mouvement du continent stable lui-même~;
par rapport à lui, Kesta s'est déplacé de $P_A = \ang{14.3}$.

\textbf{24.} $I = \ang{28.18}$ et \ang{46.16} donnent $\lambda = \ang{15.00}$ et
\ang{27.50}~: une demi-largeur de \ang{6.25}, et un écart-type de
$6.25/1.96 = \ang{3.19}$, soit \qty{355}{km}. Depuis la formation A, Kesta a dérivé
vers le nord de
\textbf{\qty[text-series-to-math=true]{2918 \pm 355}{km}}.
\end{solution}
